% Fuentes de la calibracion:
% notebooks/16_2_zenteno_recalibration_legacy_biomass_weight.executed.ipynb
% notebooks/16_3_zenteno_gf_crecimiento_azucar_disponible.executed.ipynb
% notebooks/17_2_zenteno_single_sugar_legacy_biomass_weight.executed.ipynb
% Las ecuaciones omiten las protecciones numericas de 1e-8 en denominadores.
\renewcommand{\tableshortname}{Tabla}
\chapter{Estimaci\'on de biomasa y nitr\'ogeno}
\label{cap:estimacion}

\section{Modelo cin\'etico y din\'amico del biorreactor}
\label{sec:modelo-biorreactor}

Para estimar la biomasa viable y el nitr\'ogeno asimilable por la
levadura se adaptan los balances de Zenteno y colaboradores~\cite{zenteno2010}
a un reactor de un solo compartimento. El modelo G/F considera la
glucosa y la fructosa por separado. El modelo S las representa mediante
su suma. En ambos modelos, \(X_v\) es la concentraci\'on de biomasa
viable, \(N\) la de nitr\'ogeno asimilable, \(G\) la de glucosa, \(F\)
la de fructosa, \(S=G+F\) la de az\'ucares fermentables y \(E\) la de
etanol. Las concentraciones se expresan en \(\mathrm{g/L}\), el tiempo
\(t\) en horas y la temperatura del mosto \(T\) en kelvin. Se utiliza
el perfil experimental de temperatura como entrada del modelo y se
considera que el mosto est\'a bien mezclado.

\subsection{Crecimiento, muerte y adici\'on de nitr\'ogeno}

El nitr\'ogeno asimilable limita la tasa espec\'ifica de crecimiento
\(\mu\). El crecimiento cesa cuando se agotan los az\'ucares
fermentables. La tasa espec\'ifica de muerte \(k_D\) depende de la
temperatura y del etanol.
\begin{align}
  \mu &=
    \begin{cases}
      \mu_{\max}(T)\dfrac{N}{N+K_N(T)}, & S>0,\\
      0, & S=0,
    \end{cases}
    \label{eq:c4-crecimiento}\\
  k_D &=
    \begin{cases}
      k_{D0}\exp\!\left[C_{de}E+
        \dfrac{E_{td}(T-305{,}65)}{305{,}65\,R\,T}\right], & T\ge T_D,\\
      0, & T<T_D,
    \end{cases}
    \label{eq:c4-muerte}\\
  T_D &= 315{,}89-0{,}1279 E+0{,}0049 E^2-0{,}0001 E^3 .
    \label{eq:c4-temperatura-muerte}
\end{align}

Entre adiciones de nitr\'ogeno, los balances de biomasa y
nitr\'ogeno asimilable son
\begin{equation}
  \frac{dX_v}{dt}=(\mu-k_D)X_v,\qquad
  \frac{dN}{dt}=-\frac{\mu}{Y_{X/N}}X_v .
  \label{eq:c4-balances-x-n}
\end{equation}
El coeficiente \(Y_{X/N}\) es el rendimiento de biomasa por unidad de
nitr\'ogeno consumido. Una adici\'on en \(t_i\) se representa por el
salto \(N(t_i^+)=N(t_i^-)+\Delta N_i\), donde \(\Delta N_i\) es la
concentraci\'on de nitr\'ogeno aportada. Si una muestra coincide con
una adici\'on, se utiliza el estado inmediatamente anterior al salto.

\subsection{Modelo G/F}

Las tasas espec\'ificas de producci\'on de etanol a partir de glucosa y
fructosa son
\begin{align}
  \beta_G &=
    \beta_{G\max}(T)\frac{G}{G+K_G(T)}
    \frac{K_{IE}(T)}{E+K_{IE}(T)}, \label{eq:c4-beta-g}\\
  \beta_F &=
    \beta_{F\max}(T)\frac{F}{F+K_F(T)}
    \frac{K_{IG}(T)}{G+K_{IG}(T)}
    \frac{K_{IE}(T)}{E+K_{IE}(T)} . \label{eq:c4-beta-f}
\end{align}
El etanol inhibe el consumo de ambos az\'ucares. La glucosa, adem\'as,
inhibe el consumo de fructosa. Los balances de az\'ucar incluyen el
consumo para el crecimiento, la producci\'on de etanol y el
mantenimiento celular.
\begin{align}
  \frac{dG}{dt} &=
    -X_v\left[\frac{\mu}{Y_{X/G}}+\frac{\beta_G}{Y_{E/G}}
               +m(T)\frac{G}{G+F+K_m}\right], \qquad G>0,
    \label{eq:c4-balance-g}\\
  \frac{dF}{dt} &=
    -X_v\left[\frac{\mu}{Y_{X/F}}+\frac{\beta_F}{Y_{E/F}}
               +m(T)\frac{F}{G+F+K_m}\right], \qquad F>0,
    \label{eq:c4-balance-f}\\
  \frac{dE}{dt} &= (\beta_G+\beta_F)X_v .
    \label{eq:c4-balance-e-gf}
\end{align}
El t\'ermino \(K_m=1\,\mathrm{g/L}\) aten\'ua el consumo de az\'ucar
por mantenimiento cerca del agotamiento. Al agotarse un az\'ucar, su
concentraci\'on queda en cero. El crecimiento y la captaci\'on de
nitr\'ogeno pueden continuar mientras quede glucosa o fructosa.

\subsection{Modelo S}

El modelo S conserva los balances de \(X_v\) y \(N\) de la
ecuaci\'on~\eqref{eq:c4-balances-x-n} y utiliza una sola tasa
espec\'ifica de producci\'on de etanol.
\begin{align}
  \beta_S &=
    \beta_{S\max}(T)\frac{S}{S+K_S(T)}
    \frac{K_{IE}(T)}{E+K_{IE}(T)}, \label{eq:c4-beta-s}\\
  \frac{dS}{dt} &=
    -X_v\left[\frac{\mu}{Y_{X/S}}+\frac{\beta_S}{Y_{E/S}}
               +m(T)\frac{S}{S+K_m}\right], \label{eq:c4-balance-s}\\
  \frac{dE}{dt} &= \beta_S X_v .
    \label{eq:c4-balance-e-s}
\end{align}
El rendimiento \(Y_{X/S}\) se obtiene al sumar los consumos de glucosa
y fructosa asociados al crecimiento cuando ambos az\'ucares est\'an
presentes.
\begin{equation}
  \frac{1}{Y_{X/S}}=\frac{1}{Y_{X/G}}+\frac{1}{Y_{X/F}} .
  \label{eq:c4-rendimiento-s}
\end{equation}

\subsection{Dependencia con la temperatura}

Se emplea la forma de Arrhenius utilizada por Zenteno y
colaboradores~\cite{zenteno2010}. Para cada par\'ametro \(p\) de la
Tabla~\ref{tab:c4-dependencia-termica}, \(p_0\) es su valor a la
temperatura de referencia \(T^*\) y \(E_a\) es la energ\'ia de activaci\'on.
\begin{equation}
  p(T)=p_0\exp\!\left[\frac{E_a(T-T^*)}{T^*RT}\right].
  \label{eq:c4-arrhenius}
\end{equation}
La constante de los gases es \(R=8{,}314\,\mathrm{kJ/(kmol\,K)}\).
Para la mortalidad se utiliza \(E_{td}=130\,000\,\mathrm{kJ/kmol}\) y
\(C_{de}=0{,}0415\,\mathrm{L/g}\) de etanol, con el umbral \(T_D\)
de la ecuaci\'on~\eqref{eq:c4-temperatura-muerte}.

\begin{table}[htbp]
  \centering
  \caption{Constantes para la dependencia t\'ermica de la cin\'etica.}
  \label{tab:c4-dependencia-termica}
  \begin{tabular}{lcc}
    \hline
    Par\'ametro \(p\) & \(T^*\) [K] & \(E_a\) [kJ/kmol] \\
    \hline
    \(\mu_{\max}\) & \(300\) & \(59\,453\) \\
    \(\beta_{G\max},\ \beta_{F\max},\ \beta_{S\max}\)
      & \(296{,}15\) & \(11\,000\) \\
    \(K_N,\ K_G,\ K_F,\ K_S,\ K_{IG},\ K_{IE}\)
      & \(293{,}15\) & \(46\,055\) \\
    \(m\) & \(293{,}3\) & \(37\,681\) \\
    \hline
  \end{tabular}
\end{table}

\subsection{Producci\'on y liberaci\'on de di\'oxido de carbono}

La producci\'on de di\'oxido de carbono incluye un aporte del
crecimiento y otro de la formaci\'on de etanol, como en Zenteno y
colaboradores~\cite{zenteno2010}. Se aplica un factor de activaci\'on \(A\)
en funci\'on del az\'ucar consumido \(\Delta S=S(0)-S(t)\).
\begin{equation}
  q_{\mathrm{prod}}=A X_v\left(\frac{\mu}{Y_{X/\mathrm{CO_2}}}
                         +\frac{\beta}{Y_{E/\mathrm{CO_2}}}\right),
  \qquad \beta=
  \begin{cases}
    \beta_G+\beta_F, & \text{modelo G/F},\\
    \beta_S, & \text{modelo S}.
  \end{cases}
  \label{eq:c4-co2-produccion}
\end{equation}
Los rendimientos \(Y_{X/\mathrm{CO_2}}\) y \(Y_{E/\mathrm{CO_2}}\)
son \(12{,}72\,\mathrm{g\,X/g\,CO_2}\) y
\(1{,}05\,\mathrm{g\,E/g\,CO_2}\), respectivamente. \(A\) vale cero
hasta un consumo de \(5\,\mathrm{g/L}\) y uno desde
\(10\,\mathrm{g/L}\). Entre ambos valores,
\(A=z^2(3-2z)\) con \(z=(\Delta S-5)/5\), al expresar
\(\Delta S\) en \(\mathrm{g/L}\).

El balance de di\'oxido de carbono disuelto \(C_d\) limita la
absorci\'on a la producci\'on disponible y permite la desorci\'on
cuando \(C_d>C_{\mathrm{sat}}\). Se parte de \(C_d(0)=0\).
\begin{align}
  \frac{dC_d}{dt}
    &= \min\!\left(q_{\mathrm{prod}},\,k_La\max(C_{\mathrm{sat}}-C_d,0)\right)
       -k_La\max(C_d-C_{\mathrm{sat}},0), \label{eq:c4-co2-disuelto}\\
  q_{\mathrm{gas}} &= q_{\mathrm{prod}}-\frac{dC_d}{dt}, \qquad
  q_{\mathrm{mod}} = g_{\mathrm{CO_2}}\,q_{\mathrm{gas}} .
    \label{eq:c4-co2-flujo}
\end{align}
Las tasas \(q\) se expresan por unidad de volumen de mosto, en
\(\mathrm{g\,CO_2/(L\,h)}\). \(k_La\) es el coeficiente efectivo de
transferencia y \(g_{\mathrm{CO_2}}\) es la ganancia adimensional de la
salida del modelo.

La concentraci\'on de saturaci\'on se calcula con la correlaci\'on de
Lonvaud-Funel y Matsumoto, reproducida por Liger-Belair y
colaboradores~\cite{ligerbelair2023}.
\begin{equation}
  C_{\mathrm{sat}}
  =3{,}295\,p_{\mathrm{CO_2}}\,
    \exp\!\left[(7\times10^{-5}a-0{,}02905)T_C
                  -0{,}0179a-0{,}00111S\right].
  \label{eq:c4-solubilidad}
\end{equation}
En esta correlaci\'on \(C_{\mathrm{sat}}\) queda en \(\mathrm{g/L}\)
al usar \(T_C=T-273{,}15\) en grados Celsius, \(a=E/7{,}8924\)
en porcentaje de etanol en volumen, \(S\) en \(\mathrm{g/L}\) y
\(p_{\mathrm{CO_2}}\) en bar. Se adopt\'o una presi\'on parcial de
\(1\,\mathrm{bar}\) y se aplic\'o la correlaci\'on de vino al mosto
en fermentaci\'on.

\subsection{Calibraci\'on}

Se calibraron ambos modelos con LAB004--LAB009 mediante m\'inimos
cuadrados ponderados, con cotas y dos puntos de partida. Se utilizaron
las observaciones posteriores al inicio y se conservaron las condiciones
iniciales corregidas, las temperaturas y las adiciones de nitr\'ogeno.
Para S se sumaron las mediciones simult\'aneas de glucosa y fructosa.

Cada residuo se dividi\'o por la escala de su variable. Para la
biomasa se utiliz\'o
\(\sigma_X=\max(0{,}06,\ 0{,}08|X_{v,\mathrm{obs}}|)\),
con la concentraci\'on y su escala en \(\mathrm{g/L}\). La escala de
az\'ucar total se obtuvo como \(\sqrt{\sigma_G^2+\sigma_F^2}\).
Para el nitr\'ogeno se utiliz\'o el mayor valor entre
\(0{,}012\,\mathrm{g\,N/L}\) y el \(8\) por ciento de la concentraci\'on
observada. Para glucosa y fructosa se utiliz\'o el mayor valor entre
\(2{,}5\,\mathrm{g/L}\) y el \(2{,}5\) por ciento de la concentraci\'on
observada. Para el etanol se utiliz\'o el mayor valor entre
\(2\,\mathrm{g/L}\) y el \(2{,}5\) por ciento de la concentraci\'on
observada.
En las muestras empleadas para el ajuste, las lecturas de
nitr\'ogeno asimilable de hasta \(20\,\mathrm{mg\,N/L}\) se
llevaron a cero y se excluyeron las observaciones tomadas m\'as de
12 horas despu\'es del agotamiento de nitr\'ogeno. En LAB006 se
fij\'o en cero la observaci\'on de las 122 horas y se aplic\'o
el corte despu\'es de las 134 horas. Se conserv\'o la
concentraci\'on inicial de nitr\'ogeno de cada lote.

Para G/F se conservaron los ocho par\'ametros obtenidos en el ajuste
previo, en el que el crecimiento se suspend\'ia al agotarse el
primero de los dos az\'ucares. Con esos valores se simul\'o la
regla de crecimiento de la ecuaci\'on~\eqref{eq:c4-crecimiento}.
El promedio de los cuadrados de los residuos ponderados en los
lotes de calibraci\'on fue \(4{,}7155\) para ambas reglas. En el
modelo S se ajustaron cinco par\'ametros y se obtuvo un promedio
de \(5{,}4510\) para las variables \(X_v\), \(N\), \(S\) y \(E\).

Las Tablas~\ref{tab:c4-parametros-gf} y \ref{tab:c4-parametros-s}
presentan los valores utilizados. Se fij\'o \(K_{N0}\) para obtener
\(K_N(296{,}15\,\mathrm{K})=0{,}010\,\mathrm{g\,N/L}\), seg\'un
Zenteno y colaboradores~\cite{zenteno2010}. El valor de \(K_{IE0}\) para S
procede de una calibraci\'on anterior de G/F. Los par\'ametros de
la salida de di\'oxido de carbono se conservaron de los ajustes
previos en LAB019--LAB021.

\begin{table}[htbp]
  \centering
  \caption{Par\'ametros del modelo G/F.}
  \label{tab:c4-parametros-gf}
  \small
  \setlength{\tabcolsep}{3pt}
  \begin{tabular}{@{}l>{\raggedright\arraybackslash}p{4.3cm}clr@{}}
    \hline
    S\'imbolo & Descripci\'on & Unidad & Origen & Valor \\
    \hline
    \(\mu_0\) & Tasa espec\'ifica m\'axima de crecimiento & \(\mathrm{h^{-1}}\)
      & Ajustado & \(0{,}086722\) \\
    \(\beta_{G0}\) & Tasa espec\'ifica m\'axima de producci\'on de etanol a partir de glucosa
      & \(\mathrm{g\,E/(g\,X\,h)}\) & Ajustado & \(0{,}428291\) \\
    \(\beta_{F0}\) & Tasa espec\'ifica m\'axima de producci\'on de etanol a partir de fructosa
      & \(\mathrm{g\,E/(g\,X\,h)}\) & Ajustado & \(0{,}593521\) \\
    \(K_{N0}\) & Constante de saturaci\'on de nitr\'ogeno & \(\mathrm{g\,N/L}\)
      & Fijo & \(0{,}008258\) \\
    \(K_{G0}\) & Constante de saturaci\'on de glucosa & \(\mathrm{g\,G/L}\)
      & Fijo & \(4{,}400775\) \\
    \(K_{F0}\) & Constante de saturaci\'on de fructosa & \(\mathrm{g\,F/L}\)
      & Fijo & \(19{,}296084\) \\
    \(K_{IG0}\) & Constante de inhibici\'on por glucosa & \(\mathrm{g\,G/L}\)
      & Ajustado & \(80{,}490799\) \\
    \(K_{IE0}\) & Constante de inhibici\'on por etanol & \(\mathrm{g\,E/L}\)
      & Ajustado & \(168{,}120412\) \\
    \(k_{D0}\) & Tasa espec\'ifica de muerte & \(\mathrm{h^{-1}}\)
      & Fijo & \(0{,}000440\) \\
    \(Y_{X/N}\) & Rendimiento de biomasa por nitr\'ogeno & \(\mathrm{g\,X/g\,N}\)
      & Ajustado & \(5{,}272368\) \\
    \(Y_{X/G}\) & Rendimiento de biomasa por glucosa & \(\mathrm{g\,X/g\,G}\)
      & Fijo & \(0{,}966094\) \\
    \(Y_{X/F}\) & Rendimiento de biomasa por fructosa & \(\mathrm{g\,X/g\,F}\)
      & Fijo & \(1{,}109896\) \\
    \(Y_{E/G}\) & Rendimiento de etanol por glucosa & \(\mathrm{g\,E/g\,G}\)
      & Ajustado & \(0{,}459547\) \\
    \(Y_{E/F}\) & Rendimiento de etanol por fructosa & \(\mathrm{g\,E/g\,F}\)
      & Ajustado & \(0{,}510598\) \\
    \(m_0\) & Coeficiente de mantenimiento & \(\mathrm{g\,S/(g\,X\,h)}\)
      & Fijo & \(0{,}018339\) \\
    \(k_La\) & Transferencia de di\'oxido de carbono
      & \(\mathrm{h^{-1}}\) & Ajuste previo & \(0{,}150080\) \\
    \(g_{\mathrm{CO_2}}\) & Ganancia de la salida de di\'oxido de carbono
      & Adimensional & Ajuste previo & \(0{,}646064\) \\
    \hline
  \end{tabular}
\end{table}

\begin{table}[htbp]
  \centering
  \caption{Par\'ametros del modelo S.}
  \label{tab:c4-parametros-s}
  \small
  \setlength{\tabcolsep}{3pt}
  \begin{tabular}{@{}l>{\raggedright\arraybackslash}p{4.3cm}clr@{}}
    \hline
    S\'imbolo & Descripci\'on & Unidad & Origen & Valor \\
    \hline
    \(\mu_0\) & Tasa espec\'ifica m\'axima de crecimiento & \(\mathrm{h^{-1}}\)
      & Ajustado & \(0{,}085183\) \\
    \(\beta_{S0}\) & Tasa espec\'ifica m\'axima de producci\'on de etanol a partir de az\'ucar total
      & \(\mathrm{g\,E/(g\,X\,h)}\) & Ajustado & \(0{,}899043\) \\
    \(K_{N0}\) & Constante de saturaci\'on de nitr\'ogeno & \(\mathrm{g\,N/L}\)
      & Fijo & \(0{,}008258\) \\
    \(K_{S0}\) & Constante de saturaci\'on de az\'ucar total & \(\mathrm{g\,S/L}\)
      & Ajustado & \(18{,}439784\) \\
    \(K_{IE0}\) & Constante de inhibici\'on por etanol & \(\mathrm{g\,E/L}\)
      & Fijo & \(96{,}265166\) \\
    \(k_{D0}\) & Tasa espec\'ifica de muerte & \(\mathrm{h^{-1}}\)
      & Fijo & \(0{,}000440\) \\
    \(Y_{X/N}\) & Rendimiento de biomasa por nitr\'ogeno & \(\mathrm{g\,X/g\,N}\)
      & Ajustado & \(5{,}316422\) \\
    \(Y_{X/S}\) & Rendimiento de biomasa por az\'ucar total & \(\mathrm{g\,X/g\,S}\)
      & Calculado & \(0{,}516507\) \\
    \(Y_{E/S}\) & Rendimiento de etanol por az\'ucar total & \(\mathrm{g\,E/g\,S}\)
      & Ajustado & \(0{,}484283\) \\
    \(m_0\) & Coeficiente de mantenimiento & \(\mathrm{g\,S/(g\,X\,h)}\)
      & Fijo & \(0{,}018339\) \\
    \(k_La\) & Transferencia de di\'oxido de carbono
      & \(\mathrm{h^{-1}}\) & Ajuste previo & \(0{,}137966\) \\
    \(g_{\mathrm{CO_2}}\) & Ganancia de la salida de di\'oxido de carbono
      & Adimensional & Ajuste previo & \(0{,}658678\) \\
    \hline
  \end{tabular}
\end{table}
